% Generated draft: requires mathematical and source review.
\documentclass[11pt]{article}
\usepackage[margin=25mm]{geometry}
\usepackage{fontspec}
\setmainfont{Latin Modern Roman}
\usepackage{amsmath,amssymb,mathtools,unicode-math}
\setmathfont{Latin Modern Math}
\usepackage{enumitem,needspace,xurl,hyperref}
\setlength{\parindent}{0pt}
\setlength{\parskip}{5pt}
\setlength{\emergencystretch}{3em}
\newcommand{\N}{\mathbb{N}}
\newcommand{\Z}{\mathbb{Z}}
\newcommand{\Q}{\mathbb{Q}}
\newcommand{\R}{\mathbb{R}}
\newcommand{\C}{\mathbb{C}}
\newcommand{\F}{\mathbb{F}}
\newcommand{\A}{\mathbb{A}}
\newcommand{\PP}{\mathbb{P}}
\newcommand{\E}{\mathbb{E}}
\newcommand{\eps}{\varepsilon}
\newcommand{\dd}{\,\mathrm{d}}
\begin{document}
\section*{UnsolvedMath ID: 30001788}
\textbf{Source problem code:} catalogue entry
\subsection{Definitions and mathematical statement}
Let \(X\) be a finite connected CW complex, and let \(\widetilde X\to X\) be the supplied \(G\)-cover. Interpreting \(G\)-cover in the usual regular-covering sense, \(G\) acts freely by deck transformations and \(G\backslash\widetilde X=X\). Let \((G_i)_{i\in\N}\) satisfy
\[
G_i\triangleleft G,\qquad n_i=[G:G_i]<\infty,\qquad \bigcap_{i\in\N}G_i=\{1\},
\]
and put \(X_i=G_i\backslash\widetilde X\). The displayed hypotheses do not require \(G_{i+1}\subseteq G_i\); whether the source intends this additional residual-chain condition is unresolved.

For each integer \(q\geq0\), define the rational Betti number and integral homology torsion growth quantity by
\[
b_q(X_i;\Q)=\dim_{\Q}H_q(X_i;\Q),\qquad
h_q(X_i)=\log\left|\operatorname{Tor}_{\Z}H_q(X_i;\Z)\right|.
\]
Here \(\operatorname{Tor}_{\Z}H_q(X_i;\Z)\) denotes the finite subgroup of elements of finite order, and the trivial subgroup has order \(1\). Rational coefficients are an explicit convention for these auxiliary Betti numbers; the supplied statement does not specify the coefficients of its Betti numbers. Write
\[
b_q^{(2)}(\widetilde X;G)
=\dim_{\mathcal N(G)}\overline H_q^{(2)}(\widetilde X)
\]
for the \(L^2\)-Betti number, where \(\mathcal N(G)\) is the group von Neumann algebra and \(\overline H_q^{(2)}\) is reduced \(L^2\)-homology. Write \(\rho^{(2)}(\widetilde X;G)\) for \(L^2\)-torsion only when the hypotheses required by the chosen definition hold.

\textbf{Source approximation question, with unresolved qualifications.} For the coverings above, the record asks whether normalized torsion invariants of \(X_i\) converge to the corresponding \(L^2\)-torsion of the \(G\)-cover. Its intended equality can only be represented schematically as
\[
\lim_{i\to\infty}\frac{\mathcal T(X_i)}{n_i}
=\rho^{(2)}(\widetilde X;G),
\]
where \(\mathcal T\) is a placeholder for the source's unspecified torsion invariant, not a defined invariant. The record does not determine its normalization, sign, or whether logarithms are intended. It also does not supply sufficient hypotheses ensuring that the right-hand side is defined. Consequently this display is an unresolved formulation template, not a complete conjecture.

The record also requests an analogous approximation for analytic torsion of closed Riemannian manifolds and an approximation involving integral homology torsion under the condition
\[
b_q^{(2)}(\widetilde X;G)=0\qquad\text{for every }q\geq0.
\]
For integral homology, the quantities \(h_q(X_i)/n_i\) above are well defined, but the record does not say whether the desired approximation concerns individual degrees or an alternating sum, or provide its exact equality with \(L^2\)-torsion. For analytic torsion, the finite-cover metrics, coefficient system, treatment of nonzero cohomology, and torsion convention need specification.

Finally, for an aspherical closed manifold \(M\) of dimension \(d\), meaning that its universal cover is contractible, the record asks for limiting normalized Betti numbers and integral homology torsion growth in finite covers. It expressly singles out exceptional middle-dimensional cases governed by \(d\) and the Euler characteristic \(\chi(M)\). Neither the exact formulas nor their exceptional degrees are supplied, and the covering system for this part is not explicitly identified. These missing conclusions cannot be reconstructed as a precise mathematical assertion from this record alone.
\subsection{Short English statement}
The source asks whether torsion invariants of finite covers, divided by their covering degrees, approximate \(L^2\)-torsion, with related analytic and integral homology versions. It also asks for Betti number and torsion growth limits for closed aspherical manifolds, including exceptional middle dimensions. The supplied record omits essential conventions, existence hypotheses, and the actual middle-dimensional formulas, so a complete precise conjecture remains unresolved.
\subsection{Sources}
\begin{itemize}
\item Supplied local registry record: permanent UnsolvedMath ID \(30001788\), source problem code \texttt{catalogue entry}, title ``Approximation Conjectures for \(L^2\)-Torsion and Homology.'' Project registry: \texttt{data/unsolvedmath/problems.json}; additional-record catalog: \texttt{catalog/problems.json}. This draft uses the supplied record without independently inspecting those files.
\item Supplied report attribution: \emph{Finite-dimensional Approximations of Discrete Groups} (2011), Oberwolfach Reports, pp.\ 1429--1467, DOI \texttt{10.4171/owr/2011/26}. Original-conjecture link: \url{https://doi.org/10.4171/owr/2011/26}. The record attributes the problem to workshop open-problem discussions, gives proposed year 2011, and names no individual proposer or precise problem page.
\item Supplied literature source labeled ``L2 approximation survey'': \url{https://arxiv.org/abs/1501.07446}. No author or title for this reference is supplied.
\item Literature assessment checked at \texttt{2026-08-21}: the supplied status is ``partially solved.'' The supplied assessment reports many positive cases and approximation theorems while describing the full residual-chain conjecture for torsion and \(L^2\)-torsion as unresolved. This is supplied dated curation, not an independently verified literature assessment.
\item The registry labels the statement \texttt{corrected\_verified} and says it was checked against the official report. That verification claim is supplied metadata; no source was visited or independently verified for this draft.
\end{itemize}
\end{document}
